\subsection{FRM Exam 1998---Question 22 (Capital Markets): aggregation of VaR}

\textbf{Enunciado}

\emph{Considering arbitrary portfolios \(A\) and \(B\), and their combined portfolio \(C\),
which of the following relationships always holds for VARs of \(A\), \(B\), and \(C\)?}
\begin{enumerate}[label=\alph*.]
\item \( \VaR_A+\VaR_B=\VaR_C\)
\item \( \VaR_A+\VaR_B\ge \VaR_C\)
\item \( \VaR_A+\VaR_B\le \VaR_C\)
\item \hl{None of the above}
\end{enumerate}

\subsubsection*{Solución}

Para portafolios arbitrarios, ninguna de (a)--(c) es una identidad universal. En particular:
\begin{itemize}
\item VaR \textbf{puede} violar subaditividad \((b)\) (existen contraejemplos discretos).
\item La diversificación hace que muchas veces \(\VaR_C<\VaR_A+\VaR_B\), por lo que \((c)\) tampoco es universal.
\item La igualdad \((a)\) solo ocurre en casos especiales (p.\,ej. comonotonía).
\end{itemize}

Sea \(X_A\) la pérdida del portafolio \(A\), \(X_B\) la pérdida del portafolio \(B\), y
\(X_C=X_A+X_B\) la pérdida del portafolio combinado \(C\). Entonces
\[
\VaR_A=\VaR_\delta(X_A),\qquad
\VaR_B=\VaR_\delta(X_B),\qquad
\VaR_C=\VaR_\delta(X_C)=\VaR_\delta(X_A+X_B).
\]

Para portafolios arbitrarios, el VaR no satisface la propiedad de subaditividad. Es decir, no siempre ocurre que
\[
\VaR_\delta(X_A+X_B)\le \VaR_\delta(X_A)+\VaR_\delta(X_B).
\]

De hecho, en la diapositiva de \emph{Medidas de riesgo} se presenta un contraejemplo
con dos bonos idénticos e independientes, cada uno con pérdida de \(100\) con
probabilidad \(0.03\) y pérdida \(0\) con probabilidad \(0.97\). Para cada bono,
\[
\VaR_{95\%}(X_A)=\VaR_{95\%}(X_B)=0,
\]
pero para el portafolio conjunto,
\[
\VaR_{95\%}(X_A+X_B)=100.
\]
Por tanto,
\[
\VaR_A+\VaR_B=0<100=\VaR_C.
\]
Esto muestra que la afirmación
\[
\VaR_A+\VaR_B\ge \VaR_C
\]
\textbf{no} siempre es cierta; luego la opción \textbf{(b)} no es universal.

Además, la igualdad
\[
\VaR_A+\VaR_B=\VaR_C
\]
tampoco puede valer siempre, pues el mismo contraejemplo da
\[
0\neq 100.
\]
Por lo tanto, la opción \textbf{(a)} también es falsa en general.

Por otro lado, tampoco es cierto en general que
\[
\VaR_A+\VaR_B\le \VaR_C.
\]
En muchos casos la diversificación reduce el riesgo, de modo que
\[
\VaR_C<\VaR_A+\VaR_B.
\]
Así, la opción \textbf{(c)} tampoco es una identidad universal.

En conclusión, para portafolios arbitrarios no existe una relación de validez
universal entre \(\VaR_A\), \(\VaR_B\) y \(\VaR_C\) de los tipos (a), (b) o (c).
Por ello, la respuesta correcta es (d).